\section{Related Work}
In the landscape of AI-powered video analysis and education technology, several platforms have made significant strides in enhancing video accessibility and interactivity. This section reviews three satisfied projects: VideoKen, Otter.ai, and Kwizie, which share common goals with our proposed system but focus on different aspects of the video learning experience.

\subsection{VideoKen}
VideoKen is a platform designed to transform informational videos into interactive learning experiences. It leverages AI to automatically index videos, creating table-of-contents and phrase clouds that allow users to navigate to specific topics of interest instantly. One of its primary features is the generation of auto-generated chapters, where the system analyzes visual and audio tracks to segment videos into logical sections, making long content easily navigable. Additionally, VideoKen supports in-video quizzes, allowing for the embedding of assessments directly within the video timeline to test learner understanding in real-time. The platform also offers deep search capabilities, enabling users to search for spoken words or onscreen text across a video library and jump directly to the relevant moments. While VideoKen excels in navigational indexing and manual quiz embedding, it focuses primarily on enterprise learning and development rather than automated educational assessment generation.

\begin{figure}[h!]
    \centering
    \includegraphics[width=0.8\linewidth]{image/videoken_demo.png}
    \caption{VideoKen Interface showing auto-generated chapters and search.}
    \label{fig:videoken}
\end{figure}

\subsection{Otter.ai}
Otter.ai is a leading tool for voice conversations and meeting transcription, using advanced speech recognition to generate rich notes for meetings, interviews, and lectures. Its relevance to our work lies in its robust transcription and summarization capabilities. Otter provides real-time transcription with speaker identification, distinguishing between different participants in a conversation to create clear, readable records. Beyond simple transcription, the platform generates automated executive summaries of meetings, capturing key decision points and action items to save review time. It also employs keyword extraction to automatically highlight key phrases and generate a summary cloud for a quick content overview. Otter.ai demonstrates the power of converting audio to structured text data, a foundational step for any video analysis system, though it does not extend its functionality to educational quiz generation.

\begin{figure}[h!]
    \centering
    \includegraphics[width=0.8\linewidth]{image/otter_demo.png}
    \caption{Otter.ai interface displaying real-time transcription and summary keywords.}
    \label{fig:otter}
\end{figure}

\subsection{Kwizie}
Kwizie is a more recent entrant that focuses specifically on using AI to create gamified quizzes from video content, aiming to make video learning more engaging by converting passive viewing into an active process. The platform's core strength lies in AI quiz generation, where it uses Large Language Models (LLMs) to analyze video transcripts and automatically generate multiple-choice questions. To further enhance engagement, Kwizie emphasizes a gamified approach, offering a multiplayer game format suitable for classrooms and live events. It also promotes microlearning by encouraging the breakdown of long videos into shorter, quiz-interspersed segments to improve retention. Kwizie is the closest in spirit to our "Quizum" project, validating the demand for automated assessment tools. However, our system aims to integrate these features—transcription, summarization, and quiz generation—into a cohesive workflow that also emphasizes detailed video summarization alongside assessment.

\begin{figure}[h!]
    \centering
    \includegraphics[width=0.8\linewidth]{image/kwizie_demo.png}
    \caption{Kwizie platform generating gamified quizzes from video content.}
    \label{fig:kwizie}
\end{figure}

\subsection{Comparison}
Table \ref{tab:related_work_comparison} summarizes the key differences between the analyzed platforms based on core features relevant to video learning and engagement.

\begin{table}[h!]
    \centering
    \begin{tabular}{|p{3cm}|p{3.5cm}|p{3.5cm}|p{3.5cm}|}
        \hline
        \textbf{Feature} & \textbf{VideoKen} & \textbf{Otter.ai} & \textbf{Kwizie} \\
        \hline
        \textbf{Speech $\rightarrow$ Text} & Yes (search indexing) & Yes (primary feature) & Yes (infrastructure) \\
        \hline
        \textbf{Auto Summarization} & Yes (Phrase clouds \& Topics) & Yes (Executive summaries) & No (focus on segmentation) \\
        \hline
        \textbf{Auto Quiz Gen} & No (Manual embedding) & No & Yes (AI-generated) \\
        \hline
        \textbf{Interactive Quizzes} & Yes (In-video) & No & Yes (Gamified) \\
        \hline
        \textbf{Chatbot / AI Q\&A} & Search-based & Yes (Otter Chat) & No \\
        \hline
        \textbf{Target Audience} & Enterprise L\&D & Business/Meetings & Education/Events \\
        \hline
        \textbf{Differentiation} & Navigational Search & Meeting Productivity & Gamified Assessment \\
        \hline
    \end{tabular}
    \caption{Comparison of related work projects.}
    \label{tab:related_work_comparison}
\end{table}
